\section{Consequences and open questions}\label{sec:synthesis}
The common question is whether locally admissible representations can be
made simultaneous. For envelope comparisons, a simultaneous representation
is one probability law with all the prescribed means. Positive coefficients
permit common upper attainment and reduce the comparison to lower-envelope
rounding; signs instead lead to cut ranges. Structural restrictions improve
the available couplings, while positive physical lower bounds change the
coefficients and may change the relevant incidence structure. In spatial
certification, a pseudoexpectation can satisfy a node's prescribed moments
and polynomial constraints without being a distribution on feasible
points. Counting the regions that preserve it needs both a degree
calculation and a witness-counting argument. The refinements above (the
balanced orientation constant $\beta_N$, general-radix attainment, and
lifts without degree loss) each keep their stated domain, oracle and
full-graph agreement hypotheses.

The open questions are quantitative or concern stronger methods.
The exact cubic constant lies between $1610000/743033$ and $31/12$;
optimality of the three-law mixture does not determine it. The optimized
harmonic upper asymptotic has no matching second-order lower term. For a
fixed physical aspect ratio the interval
$\max\{2,\rho_B\}\le C_{\rm box}(\rho_B)\le\rho_B+2$ is not closed by
the large-aspect leading constant. Exact structural constants for incidence
treewidth at least three remain open; the width-two coloring proof does
not extend merely by increasing the number of colors. On unequal-aspect
positive boxes even bipartite frequency-two exactness fails, and the
bipartite supremum remains between $3/2$ and two.

For spatial methods, arbitrary affine branching, additional coupled-cut
localizers, and convex hulls of feasible graphs are outside the proved
oracle scope. For affine branching we identify only an obstruction to our
proof, neither a fast affine tree nor a lower bound for all such trees.
